\documentclass[10pt, oneside]{article}   	% use "amsart" instead of "article" for AMSLaTeX format
\usepackage{geometry}                		% See geometry.pdf to learn the layout options. There are lots.
\geometry{letterpaper}                   		% ... or a4paper or a5paper or ... 
%\geometry{landscape}                		% Activate for rotated page geometry
%\usepackage[parfill]{parskip}    		% Activate to begin paragraphs with an empty line rather than an indent
\usepackage{graphicx}				% Use pdf, png, jpg, or eps§ with pdflatex; use eps in DVI mode
								% TeX will automatically convert eps --> pdf in pdflatex		
\usepackage{amssymb}
\usepackage{mhchem}
\usepackage{hyperref}

\title{Module 15, Part 1}
\author{Mark Peever\\ \texttt{mpeever@gmail.com}}
\date{April 17, 2026}

\begin{document}
\maketitle

\section*{Objectives}
\marginpar{0 minutes}
Refer to \href{https://drive.google.com/file/d/1_Hv3SaslAcLJrZ5UijrZkMslhQ-wCpJq/view?usp=sharing}{Module 15 Notes}.\\

By the end of this class, the students should be able to\ldots
\begin{itemize}
\item describe \textbf{Chemical Equilibrium}
\item describe \textbf{Reaction Rate}
\item describe the \textbf{Chemical Equilibrium Constant}
\end{itemize}

\section*{Welcome \& Devotion}
\marginpar{5 minutes}
\begin{itemize}
\item have one student read \href{https://www.biblegateway.com/passage/?search=Ecclesiastes\%209\%3A10–18\&version=LSB}{Ecclesiastes 9:10--18}
\end{itemize}

\section*{Module 13 Practice Problems Review}
\marginpar{10 minutes}

\begin{itemize}
\item solve p. 457 \# 4
\end{itemize}


\section*{Introduction to Chemical Equilibrium}
\marginpar{20 minutes}

\begin{itemize}
\item chemical reactions actually happen in pairs: left-to-right and right-to-left at the same time
\item \emph{equilibrium} occurs when the left-to-right and right-to-left reactions occur at the same rate
\item equilibrium takes some non-zero time to happen
\item we can define an \emph{equilibrium constant} as: $K = \frac{ [ C ]_{eq}^{c} [ D ]_{eq}^{d} } { [ A ]_{eq}^{a} [ B ]_{eq}^{b} }$
\item discuss what this looks like with graphs 
\end{itemize}

\section*{Le Chatelier's Principle}
\marginpar{40 minutes}
\begin{itemize}
\item Le Chatelier's Principle describes the effect of ``stresses" on Chemical Equilibria
\item solids and liquids aren't included in ``the principle"
\item describe effects:
\begin{enumerate}
\item if $\Delta P > 0$, the equilibrium will shift away from the side with more gas particles
\item if $\Delta T > 0$, the equilibrium will shift away from the side with more energy
\item if more of any reactant is added, the equilibrium will shift away from the side where it's a reactant
\end{enumerate}
\item we can think in terms of ``enthalpy" and ``entropy" being stresses
\end{itemize}

\section*{Questions for me}
\marginpar{5 minutes}

\section*{Assignment}
\marginpar{5 minutes}
\begin{itemize}
\item Experiment 15.1 (at home) (due 2026--04--21)
\end{itemize}



\end{document}  